% ==========================================
% BAB VII PENUTUP
% ==========================================
\cleardoublepage
\chapter{PENUTUP}
\label{chap:penutup}

Bab ini menyajikan kesimpulan yang ditarik dari hasil rancang bangun, implementasi, serta evaluasi yang telah dilakukan terhadap purwarupa sistem Equilibria. Selain itu, bab ini juga memuat saran bagi pengguna sistem dan rekomendasi pengembangan lebih lanjut untuk menyempurnakan penelitian di masa depan.

\section{Kesimpulan}
\label{sec:kesimpulan}

Berdasarkan hasil pelaksanaan tugas akhir yang mencakup proses perancangan, realisasi
mekanisme, dan implementasi sistem asesmen adaptif kolaboratif Equilibria, dapat
ditarik beberapa kesimpulan sebagai berikut:

\begin{enumerate}
    \item Telah berhasil dirancang sebuah arsitektur sistem asesmen adaptif yang
    mengintegrasikan parameter kemahiran individual berbasis algoritma \textit{Elo
    rating} dengan dimensi kontribusi sosial secara holistik. Rancangan ini
    memungkinkan sistem untuk memodelkan profil siswa tidak hanya berdasarkan
    kemampuan pengerjaan soal secara mandiri, tetapi juga keterlibatan aktif dalam
    aktivitas kolaboratif.
    
    \item Telah berhasil dikembangkan mekanisme mitigasi \textit{overpersonalization}
    yang bekerja secara proaktif melalui deteksi stagnasi kognitif (berdasarkan
    variansi perubahan skor $\Delta\theta$) dan pemilihan mitra kolaborasi yang
    heterogen menggunakan algoritma \textit{Constraint-based Re-ranking}. Mekanisme
    ini terbukti mampu memicu intervensi sosial secara otomatis untuk memecah
    isolasi pembelajaran (\textit{filter bubble}) yang dialami siswa.
    
    \item Telah berhasil diimplementasikan purwarupa sistem asesmen adaptif
    kolaboratif bernama Equilibria pada domain pembelajaran kueri SQL. Implementasi
    ini mencakup integrasi lingkungan eksekusi \textit{query} (\textit{sandbox})
    yang aman, pembaruan rating secara \textit{real-time}, serta antarmuka
    pembelajaran yang reaktif guna mendukung siklus belajar adaptif dan interaksi
    sejawat secara transparan.
\end{enumerate}

\section{Saran}
\label{sec:saran}

Berdasarkan pengalaman selama proses pengerjaan dan temuan pada tahap evaluasi, terdapat beberapa saran yang dapat dipertimbangkan oleh pengguna atau pengelola sistem:

\begin{enumerate}
    \item \textbf{Bagi Instruktur/Pengajar:} Disarankan untuk tidak membiarkan siswa menggunakan sistem adaptif individual dalam durasi yang terlalu lama tanpa adanya sesi interaksi antar-rekan, guna mencegah terjadinya frustrasi akademik akibat materi yang terlalu tersegmentasi pada satu tingkat kesulitan saja.
    \item \textbf{Bagi Siswa:} Umpan balik dari rekan sejawat sebaiknya ditinjau secara kritis sebagai bahan refleksi diri, karena proses memberikan dan menerima ulasan terbukti secara statistik membantu mempercepat penguasaan konsep yang sulit dibandingkan belajar secara mandiri.
\end{enumerate}

\section{Saran Pengembangan Selanjutnya}
\label{sec:future_works}

Penelitian ini masih memiliki keterbatasan yang dapat diperbaiki pada pengembangan selanjutnya, antara lain:

\begin{enumerate}
    \item \textbf{Rekalibrasi Parameter Stagnasi:} Ambang batas variansi ($\epsilon = 165$) yang digunakan saat ini terdeteksi terlalu sensitif (\textit{hypersensitive}). Pengembangan selanjutnya perlu mencari nilai ambang batas yang lebih optimal atau menggunakan jendela aktivitas (\textit{sliding window}) yang lebih lebar agar tidak setiap fluktuasi kecil memicu intervensi.
    \item \textbf{Peningkatan Akurasi NLP untuk Bahasa Indonesia:} Komponen \textit{Social Elo} dapat diperkuat dengan menggunakan model bahasa alami yang dilatih secara spesifik pada istilah teknis komputasi dalam bahasa Indonesia, guna meningkatkan akurasi penilaian kualitas umpan balik secara otomatis.
    \item \textbf{Integrasi Model Generatif (LLM):} Penggunaan \textit{Large Language Model} dapat dipertimbangkan untuk membantu memberikan \textit{scaffolding} awal kepada siswa saat melakukan \textit{peer review}, misalnya dengan memberikan saran poin-poin kesalahan yang perlu dikomentari pada kueri rekan.
    \item \textbf{Pengujian Skala Besar:} Melakukan pengujian pada populasi siswa yang lebih besar dan durasi yang lebih lama (satu semester penuh) untuk melihat stabilitas konvergensi skor \textit{Social Elo} serta dampak jangka panjang terhadap retensi pengetahuan siswa.
\end{enumerate}
